% Paper Summary: Proximal Policy Optimization Algorithms (Schulman et al., 2017)

\begin{papersummary}{2017}{Proximal Policy Optimization Algorithms}{John Schulman, Filip Wolski, Prafulla Dhariwal, Alec Radford, Oleg Klimov}{PPO introduced a simple yet effective policy gradient method that constrains policy updates to improve training stability, becoming the standard RL algorithm for fine-tuning large language models with human feedback.}

\summaryheading{Key Ideas}
PPO addresses the instability of policy gradient methods by constraining policy updates to remain close to the previous policy. The algorithm uses a clipped surrogate objective that prevents excessively large policy updates, maintaining training stability while achieving strong performance. PPO is simpler to implement and tune than previous methods like TRPO while matching or exceeding their performance. The algorithm balances exploration and exploitation through its conservative update strategy, making it reliable across diverse environments and tasks.

\summaryheading{Follow-on Works}
The contemporaneous preference-learning work of \hyperref[paper:christiano-2017-rlhf]{Christiano et al.\ (2017)} used TRPO and A2C; \hyperref[paper:ziegler-2019]{Ziegler et al.\ (2019)} first paired PPO with a learned reward model on language models, and \hyperref[paper:ouyang-2022]{InstructGPT (Ouyang et al., 2022)} was the first large-scale PPO-based RLHF, fine-tuning GPT-3 against human preferences under a KL penalty to the supervised policy. Constitutional AI, RLAIF, and DPO built on the RLHF paradigm. \hyperref[paper:shao-2024-grpo]{GRPO (Shao et al., 2024)} removed the learned value network entirely, estimating advantages from the mean reward of a sampled group, and became the dominant algorithm for reasoning-focused reinforcement learning in \hyperref[paper:deepseek-2025-r1]{DeepSeek-R1} and \hyperref[paper:lambert-2024-tulu3]{T\"ulu 3}.

\summaryheading{Lasting Contributions}
PPO made RLHF practical at scale, and its clipped surrogate objective survives in its successors even where the learned critic does not. Without reliable policy updates under a bounded step size, the preference-based post-training that shapes contemporary assistants would have been considerably harder to build at the required scale.

\end{papersummary}
